\chapter{Reverse-Engineering the Features}
\label{ch:featforensics}

Chapter~\ref{ch:targets} decoded the targets. This chapter turns the same
instruments on the \emph{features} --- and adds two of our own: signal
deconvolution and the lead--lag atlas. The output is a taxonomy of the 79
anonymous columns that directly drove feature engineering
(Chapter~\ref{ch:featurepatterns}), feature selection
(Chapter~\ref{ch:selection}), and architecture configuration
(Chapter~\ref{ch:attention}). Nothing in this chapter requires labels beyond
what training data provides, and everything in it transfers to any anonymized
tabular contest.

\section{Recovering the underlying signal by deconvolution}
\label{sec:deconv}

The forensic result of Chapter~\ref{ch:targets} --- responders are forward
SMAs of an underlying signal $s$ plus noise --- is an invitation. The SMA is
a \emph{linear} operator; linear operators can be inverted. If we can
reconstruct $\hat s$ from an observed responder path, we obtain the closest
thing to ground truth this dataset offers: the per-step innovation sequence
that everything else describes.

\subsection{The ridge inverse}

Stack one day's responder path into a vector $r \in \mathbb R^{T}$
($T = 968$) and write the construction as $r = A s + \eta$, where
$s \in \mathbb R^{T + w - 1}$ and $A$ is the banded forward-averaging
operator, $A_{t,\,t+j} = 1/w$ for $j = 0,\dots,w-1$. Choose the shortest
window available --- \resp{8}, $w=4$ --- because inverting a short average
loses the least resolution. Exact inversion is ill-posed: $A$ has near-null
spectral directions (the Dirichlet nulls of \S\ref{sec:fourier}) exactly
where the additive noise $\eta$ lives. The standard cure is Tikhonov
regularization:
\begin{equation}
\hat s \;=\; \big(A^\top A + \lambda I\big)^{-1} A^\top r
\;=\; M_\lambda\, r .
\label{eq:ridgedeconv}
\end{equation}
The matrix $M_\lambda$ depends only on $(T, w, \lambda)$: compute it
\emph{once} ($968^3$ flops, a third of a second) and apply it to every
(symbol, day) path as a single matrix multiplication. Deconvolving three
thousand series costs one GEMM.

\begin{lstlisting}
def deconvolution_matrix(T, w, lam):
    A = np.zeros((T, T + w - 1))
    for t in range(T):
        A[t, t:t+w] = 1.0 / w
    return np.linalg.solve(A.T @ A + lam * np.eye(T + w - 1), A.T)

M = deconvolution_matrix(968, 4, lam=0.5)   # once
S_hat = R8 @ M.T                            # all series at once: (n, 971)
\end{lstlisting}

$\lambda$ trades resolution against noise amplification; because the
spectral floor (\S\ref{sec:fourier}) told us noise exists but not its exact
size, we validated the reconstruction rather than derived $\lambda$:

\begin{jsbox}[: three sanity checks that must pass]
Any deconvolution you build should be validated against the known
construction before use. Ours (at $\lambda = 0.5$):
(1) $\corr(\resp{8}_t,\ \hat s_{t+k})$ peaks at $0.85$ at $k{=}1$, inside
the responder's own window --- the inverse is aligned and strong;
(2) \resp{6} correlates with $\hat s$ at $\approx 0.37$ averaged over
$k \in [1,20]$ (its construction window) and $\approx 0.01$ over
$k \in [-20, 0]$ --- the forward shift of \S\ref{sec:shift}, reproduced
numerically to the step;
(3) responders 0/1/2 (the $B{-}A$ differences) correlate with $\hat s$ at
$\approx 0$ --- differencing cancels the common signal, exactly as the
construction predicts. Three independent predictions of the theory, three
confirmations. Only then did we let $\hat s$ anchor everything that
follows.
\end{jsbox}

A deconvolution-free cross-check is also worth carrying: the first
difference $\Delta r^{(4)}_t \propto s_{t+4} - s_{t}$ contains the same
innovations without any matrix inversion (at the cost of mixing two time
points). Where the two disagree, trust neither.

\section{The lead--lag atlas}
\label{sec:atlas}

With $\hat s$ in hand, ask of every feature $x$ the only question that
matters: \emph{when does this feature know about the signal?} Compute, per
(symbol, day) series and then averaged,
\[
\ell_x(k) \;=\; \corr\big(x_t,\ \hat s_{t+k}\big),
\qquad k \in [-48, +48],
\]
via the same FFT machinery as \S\ref{sec:acf} (never by looping over lags).
The resulting curve --- we call the collection over all features the
\textbf{atlas} --- is the feature's timing signature. Summarize each curve
by three numbers:

\begin{itemize}
\item \textbf{alpha score}: mean of $\ell_x(k)$ over $k \in [+1, +20]$ ---
  correlation with the \emph{future} signal inside the target's window.
  This is, by construction, the part of the feature that can earn score.
  (Strict variant: $k \in [+4, +20]$, excluding lags that ridge smearing
  can contaminate --- always define a smear-proof band when your reference
  signal is itself reconstructed.)
\item \textbf{nowcast score}: mean over $k \in [-20, 0]$ --- correlation
  with the \emph{realized} path. Descriptive, not predictive.
\item \textbf{peak lag and peak value}: where the curve is largest, and how
  large.
\end{itemize}

\subsection{What the atlas found}

Running this over all 79 features on an 80-day window
(\code{artifacts/feature\_atlas/}) produced a taxonomy that no univariate
target-correlation screen could have drawn:

\begin{center}
\footnotesize
\begin{tabularx}{\textwidth}{@{}l>{\raggedright\arraybackslash}X>{\raggedright\arraybackslash}X>{\raggedright\arraybackslash}X@{}}
\toprule
group & members (exemplars) & signature & value \\
\midrule
\textbf{lead} & \feat{37}, \feat{38} & alpha $-0.067$, peak at $k{=}{+}3$;
also the largest raw target correlation ($-0.16$) & genuine foresight;
prime model inputs \\[2pt]
\textbf{pure lead} & \feat{18}, \feat{65} & alpha $-0.04$, nowcast
$\approx 0$ & foresight with no descriptive component \\[2pt]
\textbf{reversal block} & \feat{39}--\feat{60} (15 cols; strongest
\feat{60}: $-0.45$ at $k{=}{-}3$) & strong \emph{negative} nowcast at
$k \in [-4,-1]$, small \emph{positive} alpha & encode short-horizon
mean-reversion \\[2pt]
\textbf{slow SMA-like} & \feat{12}, \feat{67}, \feat{70} & triangular ACF
ramp, cutoff $\approx 160$, linearity $R^2 \approx 0.98$ & are themselves
moving averages $\Rightarrow$ deconvolvable (\S\ref{sec:featdeconv}) \\[2pt]
\textbf{market-wide} & \feat{04}--\feat{07} & $\approx 93\%$ of variance
common across symbols & market factors; natural gate inputs \\
\bottomrule
\end{tabularx}
\end{center}

Read the reversal block's signature carefully, because it is the kind of
economic structure forensics can surface: those features are \emph{high}
when the just-realized signal was \emph{negative} (nowcast $< 0$), and the
future signal is then mildly \emph{positive} (alpha $> 0$). That is a
mean-reversion signal, pre-packaged by the host. A univariate screen sees
only the modest net target correlation and ranks these columns middling;
the atlas sees \emph{why} they correlate and, more importantly, that their
predictive component has different timing from their descriptive one.

\begin{warnbox}[: the univariate screen conflates lead and nowcast]
$\corr(x_t, y_t)$ mixes two mechanisms: the feature describing the realized
window (worthless at prediction time beyond what other features already
say) and the feature leading the future window (the entire game). Two
features with identical target correlation can have opposite value. Rank
features by the alpha band of their lead--lag curve, not by raw target
correlation --- and when you later seed feature selection or attention
channels (Chapters~\ref{ch:selection}, \ref{ch:attention}), seed from the
alpha ranking.
\end{warnbox}

\subsection{The stability gate}

Structure is only usable if it persists to test time. Our rule ---
\textbf{no atlas finding is load-bearing until reproduced on a distant
window} --- gets its teeth from a pair of measurements:

\begin{itemize}
\item Atlas alpha ranking, computed independently on days 900--980 and
  days 1500--1580 (600 days apart): Spearman rank correlation
  \textbf{$+0.97$}, sign agreement 97\% among features with
  $|\text{alpha}| > 0.005$. This structure is a stationary property of the
  data-generating process. Build on it.
\item Symbol clusters from responder co-movement, same two windows:
  adjusted Rand index \textbf{$-0.07$} --- chance agreement. The clusters
  are real within a window and meaningless across windows. Do not build on
  them. (We killed a planned per-cluster-model architecture on this
  number alone; Chapter~\ref{ch:negative}.)
\end{itemize}

Same dataset, same methodology, opposite verdicts --- distinguished only by
the stability check. If you adopt one habit from this chapter, adopt this
one.

\section{Features that are themselves smoothed}
\label{sec:featdeconv}

The ACF fingerprint reader of \S\ref{sec:acf}, pointed at features rather
than targets, found that \feat{12}, \feat{67}, and \feat{70} carry the
triangular ramp with cutoff $\approx 160$: \emph{these features arrive
pre-smoothed}, roughly SMA-160 of some underlying series the host did not
otherwise expose. Two consequences:

\begin{enumerate}
\item Their \emph{recent increments} --- plausibly the informative part ---
  are attenuated by a factor of $\sim 160$ relative to their level. A model
  sees mostly stale context.
\item The remedy is the anti-smoothing toolkit: difference them at short
  lags, or ridge-deconvolve them exactly as in \S\ref{sec:deconv} with
  $w = 160$. Chapter~\ref{ch:smoothing} develops the general theory of when
  to smooth and when to un-smooth; this trio is its motivating example.
\end{enumerate}

The general forensic habit: \textbf{every tool you point at targets, point
at features too}. Targets are constructed; features are constructed; the
same fingerprints appear in both.

\section{Cross-sectional and structural profiling}

Three cheaper diagnostics complete the atlas, each a column of the summary
table our profiler emits:

\begin{itemize}
\item \textbf{Market share}: the fraction of a feature's variance explained
  by its per-(date, time) cross-symbol mean. Near 1 $\Rightarrow$ the
  feature is a market-level factor delivered to every symbol
  (\feat{04}--\feat{07} at $0.93$); near 0 $\Rightarrow$ idiosyncratic.
  Determines whether market-average engineering (Chapter
  \ref{ch:featurepatterns}) or cross-sectional ranking will do anything.
\item \textbf{Cross-day continuity}: correlation between a feature's value
  at the last timestamp of day $d$ and the first of day $d{+}1$. High
  $\Rightarrow$ a continuous process (levels, slow states); near zero
  $\Rightarrow$ the feature resets daily (intraday-scoped quantities). This
  decides whether cross-midnight rolling windows are meaningful for that
  feature or an artifact.
\item \textbf{Intraday profile share}: variance of the mean time-of-day
  curve relative to total variance. High values flag features that are
  mostly a deterministic function of the clock --- and remind you that
  \emph{time-of-day itself} is a first-class feature (it earned a top-five
  importance slot in every tree model we fit).
\end{itemize}

\section{The atlas as infrastructure}

We built the atlas as a research artifact; it became infrastructure. Its
consumers by the end of the campaign:

\begin{enumerate}
\item \textbf{Feature engineering}: the reversal block and the slow-SMA trio
  each spawned engineered families (Chapter~\ref{ch:featurepatterns}).
\item \textbf{Feature-selection seeding}: the symbolic-combination search of
  Chapter~\ref{ch:selection} improved measurably ($+0.00039 \to +0.00052$
  over baseline) when its candidate pool was seeded with atlas lead and
  reversal features instead of pure importance survivors.
\item \textbf{Attention-channel selection}: the inverted-attention model of
  Chapter~\ref{ch:attention} attends over a subset of variates; we feed it
  the atlas alpha set rather than all 130 columns.
\item \textbf{Synthetic benchmarking}: the generator of
  Chapter~\ref{ch:synthetic} plants lead features, a reversal block, slow
  SMAs, and market factors \emph{at the same indices with the same
  signatures}, so that architecture triage runs on a world with the real
  problem's anatomy.
\end{enumerate}

That is the quiet lesson of Part~II: forensics is not a warm-up act before
the real work. Done properly, it \emph{is} the real work, and the models
downstream are consumers of its output.

\begin{drillbox}
Build a miniature atlas for any panel dataset: (1) reconstruct or choose a
reference signal (a target, or its deconvolution if it is smoothed);
(2) compute FFT lead--lag curves for every feature; (3) summarize into
alpha/nowcast/peak columns; (4) recompute on a second window and report the
rank correlation. Total code: under 200 lines (ours is
\code{scripts/profile\_features.py}). Then answer: which three features
would you feed to a model restricted to three inputs? The atlas answer will
differ from the univariate-correlation answer --- and now you know which one
to trust.
\end{drillbox}
